% 问题一前半部分；由 example.tex 和 problem1_part1_preview.tex 共用。
\subsection{基于测向误差角域的交会定位模型}

针对问题一，本文将测向误差转化为关于干扰源位置的约束，通过
多角域的交集构造定位区域，再利用凸多边形的几何性质计算其直径。这
里的定位区域采用附录2所述的示向度角域交集，不额外叠加其他可能会引入
其他边界的可行域模型。

设对同一干扰源共进行 $k$ 次有效测向，第 $i$ 个检测点为
$S_i=(x_i,y_i)$，示向度为 $\theta_i$，且已知误差为 $\varepsilon=1^\circ$。
假定将检测点坐标作为已知输入，测向信息对应同一位置不变的干扰源
，且真实方向与示向度的圆周角差不超过 $\varepsilon$。以上约
束不要求误差相互独立，也不要求其服从均匀分布。本文后续统一使用
$k$ 表示检测点数，$h$ 表示定位区域凸多边形顶点数，长度单位为米。


令待定位点为 $X=(x,y)$，定义
\begin{equation}
u_i=(\cos\theta_i,\sin\theta_i),\qquad d_i=X-S_i,\qquad t=\tan\varepsilon.
\end{equation}
记二维叉积为 $a\times b=a_xb_y-a_yb_x$，则第 $i$ 次测向所限定的角域为
\begin{equation}\label{eq:p1-wedge}
W_i=\{X:u_i\times d_i-t(u_i\cdot d_i)\le0,\quad
-u_i\times d_i-t(u_i\cdot d_i)\le0\}.
\end{equation}
式\eqref{eq:p1-wedge}等价于 $|u_i\times d_i|\le t(u_i\cdot d_i)$。由于 $t>0$，两式相加即可推出 $u_i\cdot d_i\ge0$，因此该表示已经
排除了反向角域。

综合全部测向信息，定义
\begin{equation}\label{eq:p1-region}
P_k=\bigcap_{i=1}^{k}W_i=\{X\in\mathbb R^2:AX\le b\}.
\end{equation}

为所求的多边形定位区域。

由于每个角域均为两个闭半平面的交集，结合凸集的性质可知， $P_k$ 是闭凸集。满足题设误差界的干扰源真实位置 $G$ 属于每个 $W_i$，从而有 $G\in P_k$。
区域非空且有界时，可定义并取到其直径
\begin{equation}\label{eq:p1-diameter}
D(P_k)=\max_{X,Y\in P_k}\|X-Y\|_2.
\end{equation}
与题中定义相同，该量描述了区域内所有相容位置间分离距离的最大值，即为所求直径的数学表示。

\subsection{定位区域的判定与顶点构造}

将各检测点的角域约束汇总后，定位区域可统一表示为
\[
P_k=\left\{X\in\mathbb{R}^2\mid AX\le b\right\}.
\]
其中，每一行不等式表示一个半平面，定位区域就是这些半平面的公共部分。为了计算区域直径，首先判断交集的状态，再确定其边界顶点。

首先，检查是否存在同时满足全部约束的位置。若不存在，则交集为空，说明当前测向数据在给定误差范围内不能对应同一个干扰源位置，此时无法构造有效定位区域。

若交集非空，还需判断其是否有界。直观上，如果存在一个方向，使某个可行点沿该方向移动任意远后仍满足全部约束，则区域会沿该方向无限延伸。设这一方向为非零向量 $v$，则它应满足
\[
Av\le 0.
\]
这是因为，对任意可行点 $X_0$ 和任意 $\lambda\ge 0$，有
\[
A(X_0+\lambda v)
=AX_0+\lambda Av
\le b.
\]
因此，非空区域存在这样的非零方向时为无界区域，不具有有限直径；不存在时，区域有界，可以继续构造其顶点。

对于非空有界的交集，采用边界交点枚举法确定定位区域。具体步骤如下：

\begin{enumerate}
    \item \textbf{计算边界交点。}
    将每个半平面对应的边界写成直线方程，逐一计算两条非平行边界直线的交点。

    \item \textbf{筛选可行交点。}
    将交点代入全部半平面约束，只有满足所有不等式的交点才予以保留。

    \item \textbf{整理区域边界。}
    对保留点使用 Andrew 单调链算法求凸包，删除重复点和边上的冗余共线点，得到按逆时针顺序排列的顶点
    \[
    V_0,V_1,\ldots,V_{h-1}.
    \]
\end{enumerate}

这一方法能够完整确定定位区域：半平面交集的每个顶点都位于至少两条非平行约束边界的交会处，因此会出现在枚举结果中；同时，保留的交点均满足全部约束，不会引入区域外的点。由于定位区域具有凸性，这些交点的凸包恰好就是原定位区域。

每次测向提供两条非冗余角域边界，因此，对于非退化的有界定位多边形，其顶点数满足
\[
h\le 2k.
\]
实际顶点数可能更少，因为部分测向约束不会形成最终区域的边界。

若区域退化为单点，其直径直接取 $0$；若退化为线段，则取两端点间距。对于具有三个及以上顶点的凸多边形，将上述有序顶点序列输入旋转卡壳算法，继续求解区域直径。

% 可替换原5.3节；引用键 p1-toussaint1983 沿用主论文。
\subsection{定位区域直径的旋转卡壳求解}
\label{sec:p1-calipers}

\subsubsection{将区域直径转化为顶点间的最大距离}

上一节已获得定位区域的有序顶点 $V_0,V_1,\ldots,V_{h-1}$。以下将非空有界定位区域简记为 $P$。题目要求计算区域内任意两点之间距离的最大值，即
\begin{equation}
    D(P)=\max_{X,Y\in P}\|X-Y\|_2.
\end{equation}
这里的 $X$ 和 $Y$ 可以位于区域内部、边界或顶点，并不预先限定为多边形顶点。直接搜索区域内的全部点不可行，因此需要利用凸性将搜索范围缩小到有限个顶点。

由于 $P$ 是凸多边形，任意 $X,Y\in P$ 都可以表示为各顶点的加权平均：
\begin{equation}
    X=\sum_{i=0}^{h-1}\alpha_iV_i,\qquad
    Y=\sum_{j=0}^{h-1}\beta_jV_j,
\end{equation}
其中 $\alpha_i,\beta_j\ge0$，且 $\sum_i\alpha_i=\sum_j\beta_j=1$。因此，
\begin{equation}
\begin{aligned}
    \|X-Y\|_2
    &=\left\|\sum_{i,j}\alpha_i\beta_j(V_i-V_j)\right\|_2\\
    &\le\sum_{i,j}\alpha_i\beta_j\|V_i-V_j\|_2\\
    &\le\max_{i,j}\|V_i-V_j\|_2.
\end{aligned}
\end{equation}
也就是说，区域内任意两点的距离都不会超过最远顶点对的距离。另一方面，顶点本身也属于区域，故区域直径恰好为
\begin{equation}\label{eq:p1-vertexdiam}
    D(P)=\max_{0\le i,j<h}\|V_i-V_j\|_2.
\end{equation}
于是，只需在顶点之间寻找最远点对，便可得到整个连续区域的直径。穷举全部 $\binom h2$ 个顶点对需要 $O(h^2)$ 时间，下面利用旋转卡壳进一步减少搜索量。

\subsubsection{平行支撑线与候选点对}

设一条直线与凸多边形接触，且整个多边形都位于它的同一侧，则称这条直线为支撑线。用两条相对的平行支撑线夹住多边形时，两条线上分别接触的顶点构成对踵点对。旋转卡壳的基本思想是让这两条支撑线围绕多边形转动，并随转动更新接触顶点\cite{p1-toussaint1983}。

之所以可以重点搜索对踵点对，是因为最远顶点对必然具有这样的支撑关系。设最远顶点为 $A,B$，记 $D=\|B-A\|_2>0$，取单位方向向量 $u=(B-A)/D$。对任意 $X\in P$，都有
\begin{equation}\label{eq:p1-support}
    0\le(X-A)\cdot u\le D.
\end{equation}
如果左侧不等式不成立，则 $X$ 沿 $AB$ 方向位于 $A$ 的外侧，其到 $B$ 的距离将超过 $D$；如果右侧不等式不成立，则 $X$ 位于 $B$ 的外侧，其到 $A$ 的距离将超过 $D$。这两种情况都与 $A,B$ 为最远点对相矛盾。因此，整个多边形必位于过 $A,B$ 且垂直于 $AB$ 的两条直线之间，这两条直线正是相对的平行支撑线。

这一性质说明，只要完整处理旋转过程中可能实现最远距离的支撑点组合，就不会遗漏区域直径。

\subsubsection{利用三角形面积推进支撑点}

对逆时针排列的顶点，固定当前边 $V_iV_{i+1}$，定义
\begin{equation}\label{eq:p1-area}
    F_i(j)=(V_{i+1}-V_i)\times(V_j-V_i),
\end{equation}
其中下标按模 $h$ 循环计数。$F_i(j)$ 是三角形 $V_iV_{i+1}V_j$ 面积的两倍，且满足
\begin{equation}
    F_i(j)=\|V_{i+1}-V_i\|_2\,operatorname{dist}(V_j,L_i),
\end{equation}
其中 $L_i$ 是当前边所在的直线。由于边长固定，面积最大的顶点也是距离 $L_i$ 最远的顶点，即对侧支撑线的接触点。

在计算过程中，比较 $F_i(j+1)$ 与 $F_i(j)$：如果前者较大，就把对侧顶点从 $V_j$ 推进到 $V_{j+1}$，直到面积不再增加。这里最大化的是点到当前边所在直线的距离，并不直接等于区域直径；仍需计算相应顶点对的欧氏距离。

凸多边形的边方向沿边界有序变化，对应的支撑位置也按边界顺序变化。因此，当当前边移到下一条边时，对侧支撑点可以从已有位置继续向前搜索，无须回到起点。正是这种单向推进性质，避免了对每条边重新遍历全部顶点。

若 $F_i(j+1)=F_i(j)$ 且两者均达到最大值，则对侧支撑线同时经过两个相邻顶点，即与一条边重合。此时需要把两个对侧端点均纳入候选，完整处理平行边情形。

\subsubsection{算法流程与结果的可靠性}

综合上述关系，采用以下扫描步骤。下述比较首先按精确运算描述，数值实现中的容差另行讨论。
\begin{enumerate}[label=\arabic*)]
    \item 若区域为单点，返回 $D=0$；若区域为线段，返回两端点距离。对 $h\ge3$ 的凸多边形，初始化对侧指针 $j=1$ 和最大距离平方 $M=0$。
    \item 依次处理各边 $V_iV_{i+1}$。当 $F_i(j+1)>F_i(j)$ 时推进 $j$，直至到达当前边的对侧支撑位置。
    \item 分别计算 $\|V_i-V_j\|_2^2$ 和 $\|V_{i+1}-V_j\|_2^2$，与 $M$ 比较并保存较大值及其端点。
    \item 若对侧出现等面积的相邻顶点，则进一步检查 $\|V_i-V_{j+1}\|_2^2$ 和 $\|V_{i+1}-V_{j+1}\|_2^2$，更新最优结果。
    \item 扫描全部边后，输出 $D=\sqrt M$ 以及对应最远点对。
\end{enumerate}

上述扫描覆盖最远点对的原因可以结合支撑线理解：从任一最远点对对应的平行支撑线出发，在保持这两个顶点接触的条件下转动支撑线，直至其中一条线首次与多边形的一条边重合；若初始已经重合，则无需转动。此时最远点对必属于该边端点与对侧支撑端点的组合。算法依次处理每条边，并在平行边情形下检查全部端点组合，因此会检查到实现直径的点对。同时，所检查的点都属于定位区域，其距离不可能超过区域直径，故扫描所得最大值与式\eqref{eq:p1-vertexdiam}一致。

整个过程的适用性取决于交会区域非空、有界，以及顶点序列和支撑事件处理的完整性，不要求检测点等距或均匀分布，也不要求定位区域具有固定的边数。边指针遍历 $h$ 条边，对侧指针仅沿边界推进常数圈，故已有凸包上的直径扫描时间复杂度为 $O(h)$。

\subsection{计算复杂度与数值实现}

设半平面数量为 $m=O(k)$。边界交点枚举产生 $O(m^2)$ 个候选，每个候选检查 $m$ 个约束，故该阶段为 $O(m^3)$；候选点凸包计算至多为
 $O(m^2\log m)$。在凸包已知时，边指针遍历 $h$ 条边，对侧支撑指针仅沿边界推进常数圈，每次候选计算为常数时间，故直径扫描为 $O(h)$，
 仅维护指针及最优点对时额外空间为 $O(1)$。因此，本实现的交点枚举阶段主导整体代价，不能把完整定位流程的复杂度写成 $O(k)$。

为核验数值结果，另以 $O(h^2)$ 穷举求得 $D_{\mathrm{BF}}$，与卡壳值 $D_{\mathrm{RC}}$ 比较。本文算例报告两者一致时的公共值；
当前主程序同时执行两种算法，并将穷举值保存为直径字段。因此，算例结果不用于证明整个程序的线性运行时间。

第~\ref{sec:p1-calipers}~节的推导针对精确角域及精确几何运算。当前程序使用浮点交点计算和近重合、近共线顶点简化，所得结果属于数值实现；数值容差本身不是已证明的直
径误差上界。对极近乎平行、近退化的数据应提高计算精度并核验区域。若近似区域 $\widehat P$ 与真实紧区域 $P$ 已有经验证的 Hausdorff 
距离界 $d_H(P,\widehat P)\le\delta$，则由三角不等式可得
\begin{equation}
|D(P)-D(\widehat P)|\le2\delta.
\end{equation}
该式说明误差控制需要区域近似误差的独立界；本文不将程序设定容差直接代入 $\delta$。

\subsection{算例验证与定位直径分析}

\subsubsection{算法一致性与布设差异}

问题一未提供特定的检测点坐标表，因此使用程序内置的合成算例进行验证。下述算例取真实源 $G=(150,-250)$，按 $S_i=G+r_i(\cos\phi_i,\sin\phi_i)$ 
构造检测点，并令示向度为指向 $G$ 的精确方向；测向读数误差取零，但定位模型仍保留 $\pm1^\circ$ 的误差角域。故这些算例用于几何和算法验
证，不是模拟器正式测试日志，也不代表随机误差条件下的平均性能。

各站参数以 $(\phi_i,r_i)$ 表示。双站正交算例采用 $(0^\circ,1200)$ 和 $(90^\circ,1200)$；三站算例在此基础上增加 
$(210^\circ,1300)$；小交会角算例采用 $(0^\circ,1200)$ 和 $(5^\circ,1300)$；同方位算例采用 $(0^\circ,1200)$ 和 
$(0^\circ,1300)$。结果见表\ref{tab:p1-cases}。
\begin{table}[H]
\centering
\caption{不同检测布设下的角域交会结果}\label{tab:p1-cases}
\begin{tabular}{lrrrr}
\toprule
布设 & $k$ & $h$ & $D$/m & $|D_{\mathrm{RC}}-D_{\mathrm{BF}}|$/m\\
\midrule
双站正交 & 2 & 4 & 59.262511 & 0\\
增加第三站 & 3 & 6 & 59.262511 & 0\\
双站小交会角 $5^\circ$ & 2 & 4 & 1190.734431 & 0\\
双站同方位 & 2 & --- & 无有限直径 & ---\\
\bottomrule
\end{tabular}
\end{table}

新增第三站改变了区域形状与顶点数，却保留了原最远点对，故直径未变。小交会角算例的直径显著增大，反映了该布设下区域沿某一方向的延伸；
由于此例同时改变了一个检测距离，不将其视为交会角的单因素实验。同方位算例说明，有限次数观测并不自动保证角域交集有界。

此外，使用固定随机种子20260910进行既有测试，本次运行的1000个随机凸多边形中，卡壳值与穷举值的最大差为0；矩形、正24边形及点、线段等
测试亦通过。这里的0表示同一浮点环境中两种实现的结果一致，不能解释为相对于精确定位区域的绝对误差为0。随机核验用于发现实现错误，算法
的普遍适用性由第~\ref{sec:p1-calipers}~节的几何推导说明。

\subsubsection{新增检测点下的直径单调性}

为排除不同观测组之间的混杂，固定一组检测点及其已有示向度，按顺序加入新观测。由
\begin{equation}
P_{k+1}=P_k\cap W_{k+1}\subseteq P_k
\end{equation}
立即得到：只要区域保持非空有界，就有
\begin{equation}\label{eq:p1-monotone}
D(P_{k+1})\le D(P_k).
\end{equation}
该结论要求保留原观测并新增约束，不能用于比较两组彼此独立的检测布设。

本实验仍取 $G=(150,-250)$，依次使用以下 $(\phi_i/{}^\circ,r_i/\mathrm m)$：
\begin{equation*}
\begin{gathered}
(20,1200),(120,1350),(225,1100),(310,1450),\\
(55,1250),(140,1300),(275,1150),(35,1400),\\
(190,1050),(105,1420),(290,1180),(205,1330).
\end{gathered}
\end{equation*}
按前述方式生成精确示向度，取前 $k$ 个观测得到表\ref{tab:p1-nested}。各有界结果的卡壳值与穷举值均一致。

\begin{table}[H]
\centering
\caption{保留既有观测、逐站增加约束的嵌套实验}\label{tab:p1-nested}
\begin{tabular}{rrrr}
\toprule
检测点数 $k$ & 顶点数 $h$ & 直径 $D_k$/m & 较上一步降幅/\%\\
\midrule
2 & 4 & 69.309097 & ---\\
3 & 6 & 58.259381 & 15.94\\
4 & 8 & 56.960911 & 2.23\\
5 & 9 & 56.960911 & 0.00\\
6 & 9 & 53.792245 & 5.56\\
7 & 8 & 50.514546 & 6.09\\
8 & 8 & 50.514546 & 0.00\\
9 & 8 & 46.795490 & 7.36\\
10 & 8 & 46.795490 & 0.00\\
11 & 10 & 46.408654 & 0.83\\
12 & 10 & 46.408654 & 0.00\\
\bottomrule
\end{tabular}
\end{table}

检测点数由2增至12时，直径由69.309097 m降至46.408654 m，降幅约33.04\%。10次新增观测中有4次未使报告精度下的直径下降。区域顶点数并不
单调，也不总达到 $2k$ 上界；新约束可能切去已有顶点，亦可能完全冗余。

进一步地，记旧区域的最远点对集合为
\begin{equation*}
\mathcal M_k=\{(A,B)\in P_k\times P_k:\|A-B\|_2=D(P_k)\}.
\end{equation*}
对非空紧区域，由新区域最远点对的存在性可得
\begin{equation}
D(P_{k+1})=D(P_k)
\iff \mathcal M_k\cap(P_{k+1}\times P_{k+1})\ne\varnothing.
\end{equation}
因此，即使新角域切除了旧区域的一部分，只要仍保留任一对原最远点，直径就不变；严格减小直径则要求每一对原最远点至少有一个端点被排除。
这说明检测点增加带来的定位改善取决于新观测对原区域最大延伸方向的约束作用，而不能仅由检测点数量判断，也为后续检测点选址提供了几何依据。
